\documentclass[../../main.tex]{subfiles}
\begin{document}

\begin{center}
  {\fontsize{48pt}{58pt}\selectfont \textbf{2003}\normalsize\textbf{年度}}
\end{center}

\begin{table}[htb]
  \centering
  \caption{出題テーマと難易度}
  \begin{tabular}{|c|p{9cm}|c|}
    \hline
    \textbf{問題番号} & \textbf{テーマ} & \textbf{難易度} \\
    \hline
    第1問 & 二次関数の不等式と導関数の二乗の積分 & \\
    \hline
    第2問 & 複素数平面上の点が動く円弧 & \\
    \hline
    第3問 & 円錐と円柱の共通部分の体積 & \\
    \hline
    第4問 & 二次方程式の解の累乗と漸化式 & \\
    \hline
    第5問 & さいころの目の積の倍数判定と確率 & \\
    \hline
    第6問 & 円周率の下界の証明 & \\
    \hline
  \end{tabular}
\end{table}

\begin{table}[htb]
  \centering
  \caption{解答}
  \renewcommand{\arraystretch}{1.4}
  \begin{tabular}{|c|c|p{9cm}|}
    \hline
    \textbf{問題番号} & \textbf{小問} & \textbf{答え} \\
    \hline
    第1問 & -- & $I\ge\dfrac{4}{3}(11-4\sqrt{3})$ \\
    \hline
    第2問 & (1) & $\angle APB=\dfrac{\pi}{4}$ \\
    \cline{2-3}
    & (2) & $t=28$ \\
    \hline
    第3問 & (1) & $S(t)=2\theta\cos2\theta+\pi-\sin2\theta$ \\
    \cline{2-3}
    & (2) & 体積は$\pi-\dfrac{16}{9}$ \\
    \hline
    第4問 & (1) & $s_1=4$，$s_2=18$，$s_3=76$，
    $s_n=4s_{n-1}+s_{n-2}$ \\
    \cline{2-3}
    & (2) & $-1$ \\
    \cline{2-3}
    & (3) & $6$ \\
    \hline
    第5問 & (1) & $1-\left(\dfrac{5}{6}\right)^n$ \\
    \cline{2-3}
    & (2) & $1-\left(1+\dfrac{2n}{3}\right)
    \left(\dfrac{1}{2}\right)^n$ \\
    \cline{2-3}
    & (3) & $\displaystyle\lim_{n\to\infty}
    \dfrac{1}{n}\log(1-p_n)=\log\dfrac{5}{6}$ \\
    \hline
    第6問 & -- & $\pi>3.05$（証明問題） \\
    \hline
  \end{tabular}
\end{table}

\paragraph{試験形態}
6問．試験時間は要確認．

\end{document}
